% OpenStax Calculus Volume 2 -- Section 2.6 Exercises: Moments and Centers of Mass
% Exercises reproduced from OpenStax, licensed under CC BY 4.0.
% Source: https://openstax.org/books/calculus-volume-2/pages/2-6-moments-and-centers-of-mass
\documentclass{exercises}
\title{OpenStax Calculus Volume 2}
\subtitle{Section 2.6 Exercises: Moments and Centers of Mass}
\sourceurl{https://openstax.org/books/calculus-volume-2/pages/2-6-moments-and-centers-of-mass}
\author{}
\date{}

\begin{document}
\maketitle


For the following exercises, calculate the center of mass for the collection of masses given.

\begin{enumerate}[start=1]
\item $m_{1}=2$ at $x_{1}=1$ and $m_{2}=4$ at $x_{2}=2$
\item $m_{1}=1$ at $x_{1}=-1$ and $m_{2}=3$ at $x_{2}=2$
\item $m=3$ at $x=0,1,2,6$
\item Unit masses at $(x,y)=(1,0),(0,1),(1,1)$
\item $m_{1}=1$ at $(1,0)$ and $m_{2}=4$ at $(0,1)$
\item $m_{1}=1$ at $(1,0)$ and $m_{2}=3$ at $(2,2)$
\end{enumerate}

For the following exercises, compute the center of mass $\overset{-}{x}$.

\begin{enumerate}[resume]
\item $\rho=1$ for $x\in(-1,3)$
\item $\rho=x^{2}$ for $x\in(0,L)$
\item $\rho=1$ for $x\in(0,1)$ and $\rho=2$ for $x\in(1,2)$
\item $\rho=\sin x$ for $x\in(0,\pi)$
\item $\rho=\cos x$ for $x\in(0,\frac{\pi}{2})$
\item $\rho=e^{x}$ for $x\in(0,2)$
\item $\rho=x^{3}+xe^{-x}$ for $x\in(0,1)$
\item $\rho=x\sin x$ for $x\in(0,\pi)$
\item $\rho=\sqrt{x}$ for $x\in(1,4)$
\item $\rho=\ln x$ for $x\in(1,e)$
\end{enumerate}

For the following exercises, compute the center of mass $(\overset{-}{x},\overset{-}{y})$. Use symmetry to help locate the center of mass whenever possible.

\begin{enumerate}[resume]
\item $\rho=7$ in the square $0\le x\le1$, $0\le y\le1$
\item $\rho=3$ in the triangle with vertices $(0,0)$, $(a,0)$, and $(0,b)$
\item $\rho=2$ for the region bounded by $y=\cos(x)$, $y=-\cos(x)$, $x=-\frac{\pi}{2}$, and $x=\frac{\pi}{2}$
\end{enumerate}

For the following exercises, use a calculator to draw the region, then compute the center of mass $(\overset{-}{x},\overset{-}{y})$. Use symmetry to help locate the center of mass whenever possible.

\begin{enumerate}[resume]
\item \techrequired The region bounded by $y=\cos(2x)$, $x=-\frac{\pi}{4}$, and $x=\frac{\pi}{4}$
\item \techrequired The region between $y=2x^{2}$, $y=0$, $x=0$, and $x=1$
\item \techrequired The region between $y=\frac{5}{4}x^{2}$ and $y=5$
\item \techrequired Region between $y=\sqrt{x}$, $y=\ln(x)$, $x=1$, and $x=4$
\item \techrequired The region bounded by $y=0$, $\frac{x^{2}}{4}+\frac{y^{2}}{9}=1$
\item \techrequired The region bounded by $y=0$, $x=0$, and $\frac{x^{2}}{4}+\frac{y^{2}}{9}=1$
\item \techrequired The region bounded by $y=x^{2}$ and $y=x^{4}$ in the first quadrant
\end{enumerate}

For the following exercises, use the theorem of Pappus to determine the volume of the shape.

\begin{enumerate}[resume]
\item Rotating $y=mx$ around the $x$-axis between $x=0$ and $x=1$
\item Rotating $y=mx$ around the $y$-axis between $x=0$ and $x=1$
\item A general cone created by rotating a triangle with vertices $(0,0)$, $(a,0)$, and $(0,b)$ around the $y$-axis. Does your answer agree with the volume of a cone?
\item A general cylinder created by rotating a rectangle with vertices $(0,0)$, $(a,0),(0,b)$, and $(a,b)$ around the $y$-axis. Does your answer agree with the volume of a cylinder?
\item A sphere created by rotating a semicircle with radius $a$ around the $y$-axis. Does your answer agree with the volume of a sphere?
\end{enumerate}

For the following exercises, use a calculator to draw the region enclosed by the curve. Find the area $M$ and the centroid $(\overset{-}{x},\overset{-}{y})$ for the given shapes. Use symmetry to help locate the center of mass whenever possible.

\begin{enumerate}[resume]
\item \techrequired Quarter-circle: $y=\sqrt{1-x^{2}}$, $y=0$, and $x=0$
\item \techrequired Triangle: $y=x$, $y=2-x$, and $y=0$
\item \techrequired Lens: $y=x^{2}$ and $y=x$
\item \techrequired Ring: $y^{2}+x^{2}=1$ and $y^{2}+x^{2}=4$
\item \techrequired Half-ring: $y^{2}+x^{2}=1$, $y^{2}+x^{2}=4$, and $y=0$
\item Find the generalized center of mass in the sliver between $y=x^{a}$ and $y=x^{b}$ with $a>b$. Then, use the Pappus theorem to find the volume of the solid generated when revolving around the y-axis.
\item Find the generalized center of mass between $y=a^{2}-x^{2}$, $x=0$, and $y=0$. Then, use the Pappus theorem to find the volume of the solid generated when revolving around the y-axis.
\item Find the generalized center of mass between $y=b\sin(ax)$, $x=0$, and $x=\frac{\pi}{a}$. Then, use the Pappus theorem to find the volume of the solid generated when revolving around the y-axis.
\item Use the theorem of Pappus to find the volume of a torus (pictured here). Assume that a disk of radius $a$ is positioned with the left end of the circle at $x=b$, $b>0$, and is rotated around the y-axis. \\ \textit{[Figure: This figure is a torus. It has inner radius of b. Inside of the torus is a cross section that is a circle. The circle has radius a.]}
\item Find the center of mass $(\overset{-}{x},\overset{-}{y})$ for a thin wire along the semicircle $y=\sqrt{1-x^{2}}$ with unit mass. (Hint: Use the theorem of Pappus.)
\end{enumerate}

\end{document}
